\section{The complete first nonreduced contact fibre}
\label{sec:complete-contact-v163}
We determine the entire fibre, including its nilpotents, at the first
nonreduced contact. The calculation also determines the common boundary
of the two components found in Theorem~\ref{thm:nonequidimensional-v162}.
It takes place in the graph over the coefficient base: the coefficients
are part of a point throughout.

Put $B_2=\mathbb A^6$, with coordinates the coefficients of a monic
quadratic $f$ and of two polynomials $g,r$ of degree at most one. Let
$\Gamma_2$ be the reduced closure in
$B_2\times\operatorname{Hilb}_{3l+1}(\Pj^1\times\Pj^2)$ of the graphs
of $[f:g:r]$, on their common nonvanishing locus. The polarization is
$\OO(1,1)$. Write $\nu:\widetilde\Gamma_2\to\Gamma_2$ for the
normalization of the total graph, and set
\[
 o=(x^2,0,0),\qquad X_2=(\widetilde\Gamma_2)_o.
\]
The scheme $X_2$ is not defined by normalizing a fibre. In the target
use coordinates $[F:G:R]$ and put $P=[1:0:0]$. The main curve at $o$
is $C_0=\Pj^1\times\{P\}$.

\subsection{An embedded chart containing all conic tails}
Let $[s:t]$ be coordinates on the first projective line, with $x=t/s$.
Consider the affine space with coordinates $(c,e,\lambda,A,B,C)$, and put
$z=t-cs$, $H=R-\lambda G$. In the product with $\Pj^1\times\Pj^2$
form the matrix
\begin{equation}\label{eq:HB-matrix-v163}
 \mathcal M=
 \begin{pmatrix}
 H&AF+CG\\
 -G&H+BF\\
 es&-z
 \end{pmatrix}.
\end{equation}
Its minors generate the bihomogeneous ideal
\begin{equation}\label{eq:HB-ideal-v163}
 \begin{split}
 Q&=H^2+BFH+AFG+CG^2,\\
 L_1&=zG-es(H+BF),\\
 L_2&=zH+es(AF+CG).
 \end{split}
\end{equation}
The associated map to $B_2$, written on $s=1$ with $z=x-c$, is
\begin{equation}\label{eq:HB-base-v163}
 \begin{split}
 f&=z^2+e^2C,\qquad g=eBz-e^2A,\\
 r&=\lambda g-eAz-e^2BC.
 \end{split}
\end{equation}

\begin{lemma}[The determinantal family and its recovery maps]
\label{lem:HB-family-v163}
The minors in \eqref{eq:HB-ideal-v163} define a flat family of embedded
curves of Hilbert polynomial $3l+1$. It defines a morphism to $\Gamma_2$.
Near every point with $e=0$, this morphism is an open immersion. Its
image is regular and is unchanged by normalization of the total graph.
At $e=0$ the curve is the scheme-theoretic union of $C_0$ and the plane
conic $Q=0$ over $x=c$, with the attachment placed at that base
point. The intersection is the reduced point $(c,P)$, even when $Q$
is a doubled line. Constant changes of target coordinates fixing $P$
provide such charts around every plane conic through $P$.
\end{lemma}
\begin{proof}
The minors have height two on every geometric fibre. When $e\ne0$,
one eliminates $H$ locally from $L_1,L_2$; equivalently the rank-one
condition is the two-generator graph of $f,g$, with the indicated third
coordinate. At $e=0$ the ideal is
\[
 (zG,zH,Q)=(G,H)\cap(z,Q),
\]
where $Q\in(G,H)$ is nonzero. The equality follows by reducing modulo
$z$ and then using that $z$ is regular on the polynomial ring in
$F,G,H$ modulo $Q$. This gives the asserted union and intersection.
In particular the height is two also for a singular conic. The
Hilbert--Burch complex of \eqref{eq:HB-matrix-v163} is consequently
exact on every fibre. Its sheaf form is
\begin{equation}\label{eq:HB-resolution-v163}
 0\longrightarrow\OO(-1,-2)^2\longrightarrow
 \OO(-1,-1)^2\oplus\OO(0,-2)\longrightarrow
 \OO\longrightarrow\OO_{\mathcal C}\longrightarrow0.
\end{equation}
The fibrewise exactness criterion for a complex of base-flat locally
free sheaves gives base flatness of the cokernel. Computing the Euler
polynomial from this complex gives $3l+1$. This is an application of
the classical determinantal resolution, not an assertion that the
resolution theorem is new.

Substitution of \eqref{eq:HB-base-v163} into the three minors gives
zero. On the dense open
\[
             e(A^2+CB^2)\ne0
\]
the three binary quadratics $f,g,r$ are a basis, the curve is their
graph, and the image is a smooth conic. Thus the flat family gives a
morphism into the reduced graph closure. The assertion at infinity
is direct: $s=0$ implies $z=t\ne0$, so $L_1,L_2$ give $G=H=0$.
There is no additional component or choice at infinity.

We prove the open-immersion assertion by regular recovery, not by
counting parameters. At a central curve $C_0\cup Q$, the union exact
sequence and its twists give
\[
 H^1(\OO(0,2)|_{C_0\cup Q})=0,\quad h^0=5,
 \qquad H^1(\OO(1,1)|_{C_0\cup Q})=0,\quad h^0=4.
\]
Indeed $H^0(Q,\OO_Q(2))$ has dimension five and
$H^0(Q,\OO_Q(1))$ dimension three for every nonzero plane conic,
including a doubled line, by
$0\to\OO_{\Pj^2}(j-2)\to\OO_{\Pj^2}(j)\to\OO_Q(j)\to0$.
The restriction to $P$ is surjective. Combining it with the main
component proves both equalities and the surjectivity of the ambient
restriction maps. Cohomology and base change therefore give, in a
neighbourhood of this point in the integral graph $\Gamma_2$, a line
of quadratic equations and a rank-two bundle of equations of bidegree
$(1,1)$.

Normalize the coefficient of $R^2$ in the quadratic equation to one.
Its $F^2$ coefficient is zero on the dense graph, because infinity
maps to $P$, and hence is zero throughout this neighbourhood.
Completing the square in $R,G$ uniquely writes the equation as $Q$
in \eqref{eq:HB-ideal-v163}; it recovers
$\lambda,A,B,C$ regularly. The projection to $B_2$ supplies $f$,
so $c$ is minus half its linear coefficient. Normalize the two
bilinear equations by their coefficients at $tG,tH$; the relevant
minor is a unit at the central curve. The negative coefficient of
$sH$ in the equation with $tG$ coefficient one recovers $e$.

Here is an explicit verification on the dense graph, which also
checks the normalizations. Write $f=(x-c)^2+d$,
$g=eB(x-c)+u$, and $r-\lambda g=-eA(x-c)+v$. The coefficient of
$x^3$ in $Q(f,g,r-\lambda g)=0$ gives the displayed opposite linear
coefficients. When $A^2+CB^2\ne0$, the remaining coefficients give
\[
 u=-e^2A,\qquad v=-e^2BC,\qquad d=e^2C.
\]
The normalized bilinear equations are then precisely $L_1,L_2$.
All these are equalities of regular functions or sections on an
integral neighbourhood, so they extend from its dense graph open.
The recovered parameters satisfy \eqref{eq:HB-base-v163}, and the
universal curve equals \eqref{eq:HB-ideal-v163}, first densely and
then by separatedness of the Hilbert scheme. Recovery is also the
identity on the determinantal parameter space. These two morphisms
are inverse after restricting to their common open neighbourhoods.
The parameter space is smooth; its image is therefore regular and
normalization does not change it.

Finally any nonzero quadratic form vanishing at $P$ is nonzero on
some vector not proportional to $P$. Choose that vector as the $R$
direction. This is a constant target change fixing $P$, followed by
rescaling. The preceding argument applies. Such changes are changes
of chart, not an equivalence relation on Hilbert points.
\end{proof}

\subsection{The full fibre and its nilpotents}
Let $\mathbb F_2=\Pj(\OO_{\Pj^1}\oplus\OO_{\Pj^1}(2))$ be the
compactification of $\operatorname{Tot}(\OO(2))$ whose section at
infinity has self-intersection $-2$. Denote that section by $D_\infty$.
Let $\mathcal Q_P\simeq\Pj^4$ be the projective space of plane conics
through $P$, and let $D_{\rm dbl}\simeq\Pj^1\subset\mathcal Q_P$
be its locus of doubled lines through $P$.

\begin{theorem}[Complete first nonreduced fibre]
\label{thm:complete-fibre-v163}
The scheme $X_2$ has exactly two irreducible components, with reduced
structures $\mathcal Q_P\simeq\Pj^4$ and $\mathbb F_2$. Their
scheme-theoretic intersection, when each is given its reduced closed
structure, is the reduced projective line
\[
                  D_{\rm dbl}=D_\infty.
\]
The open complement of $D_\infty$ in the second component is
$J_1(\Pj^1)=\operatorname{Tot}(T\Pj^1)$, and the first component
parametrizes exactly the curves $C_0\cup Q$, $Q\in\mathcal Q_P$.
Thus the second component has dimension exactly two and the conic
component has dimension exactly four.

The fibre is connected and generically reduced on both components,
but is not reduced. In the charts of Lemma~\ref{lem:HB-family-v163}
its scheme structure is
\begin{equation}\label{eq:full-fibre-ring-v163}
 \C[\lambda,e,A,B,C]/(eA,eB,e^2C).
\end{equation}
At a doubled line with slope $\lambda_0$ the completed local ring is
obtained by replacing the polynomial ring by
$\C[[\lambda-\lambda_0,e,A,B,C]]$ in this formula. The nilradical
is generated by $eC$, has square zero, and is supported on the locus
of conics singular at $P$. The fibre is not Cohen--Macaulay at that
locus. Its reduced support is the union of the two indicated smooth
components, and the normalization of that reduced support is their
disjoint union. This last normalization is different from the
normalization of the total Hilbert graph used to define $X_2$.
\end{theorem}
\begin{proof}
First take the fibre of \eqref{eq:HB-base-v163} at $o$. Its coefficient
equations give $c=0$, $e^2C=0$, $eB=0$, $e^2A=0$ and
$e(\lambda B-A)=0$, together with a redundant constant equation.
The resulting ideal is exactly
\[
                    (c,eA,eB,e^2C).
\]
No radical has been taken. In particular the charts compute the
scheme fibre of a regular total graph and therefore of its
normalization as well.

The conic family is algebraic over all of $\mathcal Q_P$. Its ideal
is $I_{C_0}\cap I_Q$, with $Q$ placed at $x=0$; the intersection
with $C_0$ is the constant reduced point $P$. The union exact
sequence proves flatness and the Hilbert polynomial. The family
is a closed immersion into the Hilbert fibre: its inverse on its
image recovers the quadratic equation as the kernel of the ambient
restriction in degree $(0,2)$. This also gives the relative residual
construction; locally
\[
              (I_{C_0}\cap I_Q):I_{C_0}=I_Q.
\]
For the displayed ideals this follows by dividing by $G$ and $H$:
no nonzero element in the hypersurface ring of a conic is annihilated
by both $G$ and $H$. This remains true in the universal conic family
by its same presentation. Alternatively the kernel-line construction
alone proves the asserted relative closed immersion. The regular
charts of Lemma~\ref{lem:HB-family-v163} show that this family lifts
uniquely to the normalization near every one of its points; we do
not use a general lifting assertion for maps from normal schemes.

The primitive open of the other component is the established jet
chart. Write a jet as $w=w_0+w_1x$ modulo $x^2$. Near its direction
at infinity use
\[
             \lambda=w_0,\qquad e=1/w_1.
\]
In \eqref{eq:HB-ideal-v163} put $c=A=B=C=0$. This gives
\begin{equation}\label{eq:adjacency-family-v163}
         (xG-eH,\ xH,\ H^2),\qquad H=R-\lambda G.
\end{equation}
For $e\ne0$ it is the primitive thick-line curve with slope
$w=\lambda+x/e$. At $e=0$ it is
\[
               (xG,xH,H^2)=(G,H)\cap(x,H^2),
\]
the main curve together with the planar doubled line $H^2=0$ at
$x=0$. Under the other affine coordinate of the direction line,
$w_0'=1/w_0$ and $w_1'=-w_1/w_0^2$; hence
$e'=-w_0^2e$ on the boundary chart. These are exactly the transition
functions of the compactification of $T\Pj^1=\OO(2)$ described
above. Formula \eqref{eq:adjacency-family-v163} and the jet families
glue to a morphism from $\mathbb F_2$ into $X_2$. The local inverse
recovers $\lambda,e$; the original jet chart supplies the inverse
on its complement. Thus this is a closed immersion, since its
source is proper. Its infinity section is precisely $D_{\rm dbl}$.
It has self-intersection $-2$ by the transition functions. The sum
of the two reduced component ideals in \eqref{eq:full-fibre-ring-v163}
is $(e,A,B,C)$, giving the stated reduced intersection.

It remains to prove that these two families exhaust the fibre.
Their union $K$ is closed and proper. It is also open in the
underlying fibre: at a point of the conic family the local fibre
\eqref{eq:full-fibre-ring-v163} has support
\[
                    V(e)\ \cup\ V(A,B,C),
\]
which is exactly the union just constructed; at a jet point the
original open jet chart has no other branch. The scheme
$\widetilde\Gamma_2$ is integral and normal, and its projection
to the normal integral base $B_2$ is proper and birational.
Consequently all its geometric fibres are connected by Stein
factorization \cite[Lemma~37.53.6]{Stacks163}. The nonempty subset
$K$ is open and closed in $X_2$, so it is all of its support.
This is an exhaustion argument for the proper fibre, not an
inference from the existence of two families of curves.

For clarity, the local primary calculation is
\begin{equation}\label{eq:primary-fibre-v163}
 (eA,eB,e^2C)=(e)\cap(A,B,C)\cap(A,B,e^2).
\end{equation}
It is irredundant wherever all the displayed coordinate functions
vanish. It follows either by divisibility of monomials, or by
intersecting first $(e)$ with $(A,B,C)$ and then with $(A,B,e^2)$.
The two minimal primes are $(e)$ and $(A,B,C)$; $(e,A,B)$ is the
embedded associated prime when present. The radical is
$(eA,eB,eC)$, so the nilradical is generated by $eC$ and
\[
              \operatorname{Ann}(eC)=(e,A,B),\qquad (eC)^2=0.
\]
In conic coordinates the equations $A=B=0$ say exactly that $Q$
is singular at $P$. These calculations also apply after localization
and completion and identify the claimed support intrinsically.
They prove generic reducedness, nonreducedness and failure of the
Cohen--Macaulay property. The two reduced components are smooth and
are distinct; the final assertion about their normalization follows.
\end{proof}

\begin{corollary}[Explicit adjacency and a numerical consequence]
\label{cor:adjacency-v163}
Every doubled line through the attaching point is the common limit of
primitive thick-line jets and of smooth conic tails. At such a point
there is no extra component beyond the two in
Theorem~\ref{thm:complete-fibre-v163}. The Euler characteristic of the
underlying complex fibre is seven.
\end{corollary}
\begin{proof}
The first specialization is \eqref{eq:adjacency-family-v163}. For the
second, set $e=0$, $\lambda$ constant, $B=C=0$, $A=\tau$; then
$Q=H^2+\tau FG$ is nonsingular for $\tau\ne0$ and specializes to
the doubled line. These are algebraic one-parameter families.
Exhaustion was proved in the theorem. Euler characteristic is
additive for the closed union, giving $5+4-2=7$.
\end{proof}
